\documentclass[11pt]{article}

\usepackage[a4paper,margin=2.4cm]{geometry}
\usepackage{amsmath,amssymb,amsthm,mathtools}
\usepackage{array}
\usepackage{booktabs}
\usepackage{longtable}
\usepackage{enumitem}
\usepackage[colorlinks=true,linkcolor=blue,citecolor=blue,urlcolor=blue]{hyperref}
\usepackage{xcolor}
\usepackage[export]{adjustbox}
\usepackage{microtype}
\usepackage{seqsplit}
\emergencystretch=2.5em

\title{The Bottom Panel (Viz-Dock): Implementation and Workspace Sync\\
\large How Seven Tools Share One Dock, and How Each Links Back to the Blocks}
\author{Research note --- Block-MiniJava}
\date{2026-07-20}

\begin{document}
\maketitle

\begin{abstract}
The bottom panel --- \texttt{<section id="viz-dock">}
(\texttt{src/index.html:243--327}), internally called the \emph{viz-dock}
--- hosts seven tools behind one tab strip: \textbf{Problems},
\textbf{Output}, \textbf{Call-by-Structure}, \textbf{Call-by-Value},
\textbf{CESK}, \textbf{A~vs~B}, and \textbf{Rewrite}. All seven share one
shell (\texttt{src/core/ui/visualizationPanel.ts}) for opening, closing,
resizing, and tab-switching, but they fall into three genuinely different
\emph{synchronization regimes} with the main Blockly workspace: two tools
highlight the live, original blocks; three tools render an independent
\emph{copy} of a block subtree in their own private workspace and never
touch the original; and one tool is not block-based at all. This note
explains the shared shell first, then each regime with the exact functions
that implement it.
\end{abstract}

\tableofcontents

% =============================================================================
\section{The shared shell}

\texttt{src/core/ui/visualizationPanel.ts} (313 lines) owns everything
about the dock that is common to all seven tools: which one is active,
whether the dock is open, its height, and where each tool's DOM host lives.
It does not itself know how to run a machine or check types --- those are
injected as callbacks or driven by other modules that this one only calls
into.

\begin{center}
\begin{longtable}{@{}p{4.0cm}p{10.2cm}@{}}
\toprule
\textbf{Symbol} & \textbf{Role} \\
\midrule
\endhead
\texttt{VizKind} &
  \texttt{visualizationPanel.ts:15} --- the closed union
  \texttt{ReductionKind | 'machine' | 'compare' | 'subst' | 'problems' |
  'output'}, i.e. \texttt{'structure' | 'value' | 'machine' | 'compare' |
  'subst' | 'problems' | 'output'} --- one tag per tab. \\
\texttt{ALL\_KINDS} &
  \texttt{visualizationPanel.ts:18} --- the tab order as rendered:
  \texttt{problems, output, structure, value, machine, compare, subst}. \\
\texttt{TITLE} &
  \texttt{visualizationPanel.ts:19--27} --- one human-readable description
  per kind, shown in \texttt{\#viz-dock-info} (\S\ref{sec:info}). \\
\texttt{hostOf(kind)} / \texttt{tabOf(kind)} &
  \texttt{visualizationPanel.ts:62--72} --- look up
  \texttt{.viz-host[data-kind=kind]} / \texttt{.viz-tab[data-kind=kind]} in
  the DOM; every other function addresses a tab/host pair only through
  these two, so the seven tools never need seven special cases. \\
\texttt{setActive(kind)} &
  \texttt{visualizationPanel.ts:128--146} --- the single tab-switch
  function: toggles \texttt{data-active}/\texttt{hidden}/\texttt{aria-*}
  on every host and tab, calls \texttt{setStepperTabVisible} /
  \texttt{setCompareTabVisible} / \texttt{setSubstTabVisible} so the
  newly-hidden or newly-shown stepper can pause its \textbf{Play} loop or
  redraw its SVG overlay (companion note
  \texttt{cesk-machine-and-stepper.tex}, \S4.1), persists the choice to
  \texttt{localStorage}, and resizes the active tab's own Blockly
  instance. \\
\texttt{\seqsplit{openBottomTool(kind)}} &
  \texttt{visualizationPanel.ts:148--151} --- \texttt{setActive(kind)} then
  force-opens the dock; the single entry point every "jump to a specific
  bottom tool" call site uses (Problems from the status bar,
  \S\ref{sec:problems}; \texttt{Run} opening Output,
  \texttt{app.ts:340}). \\
\texttt{setVizOpen(open)} &
  \texttt{visualizationPanel.ts:157--173} --- shows/hides the whole dock
  (\texttt{dock().dataset.open}), clears "maximized" state on close,
  persists to \texttt{localStorage}, and calls the shell's
  \texttt{onLayoutChange} callback (installed by
  \texttt{initVisualizationPanel}, wired from \texttt{app.ts}'s layout
  coordinator --- see \texttt{ui-design-and-layout.tex}). \\
\texttt{initResizer()} &
  \texttt{visualizationPanel.ts:216--251} --- pointer-drag height resize
  on \texttt{\#viz-resizer}, plus \texttt{ArrowUp}/\texttt{ArrowDown}
  keyboard resize (l.245--250); persists height to \texttt{localStorage}. \\
\texttt{\seqsplit{initVisualizationPanel(layoutChange)}} &
  \texttt{visualizationPanel.ts:253--313} --- called once from
  \texttt{app.ts}; wires every tab's click/keyboard-arrow navigation
  (l.255--273), the toolbar's \textbf{Re-run}/\textbf{Arrange}/\textbf{Collapse}/
  \textbf{Maximize} buttons, restores the persisted height/open/active-tab
  state, and installs the resizer. \\
\bottomrule
\end{longtable}
\end{center}

\subsection{The info line and passive vs.\ active tools}
\label{sec:info}

\texttt{updateInfo()} (\texttt{visualizationPanel.ts:111--118}) writes
\texttt{TITLE[active]} into \texttt{\#viz-dock-info}, appending the
currently-visualized block's type for Structure/Value.
\texttt{updateToolActions()} (l.120--126) hides the \textbf{Re-run} button
for the two tools that have nothing to re-run (\texttt{problems},
\texttt{output} --- both are "passive": they mirror state pushed into them
rather than owning a run of their own) and hides \textbf{Arrange} for every
tool except Structure/Value (only they lay out a tree of copied blocks that
benefits from re-arranging).

% =============================================================================
\section{The three sync regimes}

Every artifact rendered in the bottom panel is, in principle, a candidate
for "click it to jump to the block it came from." Which workspace it jumps
\emph{to} is what separates the three regimes below.

\begin{center}
\begin{longtable}{@{}p{2.6cm}p{4.4cm}p{7.6cm}@{}}
\toprule
\textbf{Regime} & \textbf{Tabs} & \textbf{Sync target} \\
\midrule
\endhead
Live highlight &
  Problems, CESK, A~vs~B &
  The \emph{original, live} main workspace --- the same
  \texttt{Blockly.WorkspaceSvg} the user is editing. Clicking an artifact
  centers and selects the real block; a running machine continuously
  highlights its current focus block there too. \\
Independent copy &
  Call-by-Structure, Call-by-Value, Rewrite &
  A \emph{separate} Blockly workspace injected inside the dock, showing a
  freshly rendered copy of a block subtree. Clicking inside these tabs
  never touches the main workspace, because the tree on screen is not the
  same object as the blocks it was copied from. \\
Text only &
  Output &
  Not block-based at all --- a plain \texttt{<pre>} console; there is
  nothing to click-to-locate. \\
\bottomrule
\end{longtable}
\end{center}

\subsection{Live highlight: Problems}
\label{sec:problems}

\texttt{src/core/ui/typeDiagnostics.ts} (104 lines) is the smallest and
clearest example of the pattern every other live-highlight tool follows.

\begin{center}
\begin{longtable}{@{}p{4.0cm}p{10.2cm}@{}}
\toprule
\textbf{Function} & \textbf{Role} \\
\midrule
\endhead
\texttt{\seqsplit{refreshTypeDiagnostics(workspace)}} &
  \texttt{typeDiagnostics.ts:99--104} --- the one exported entry point:
  runs \texttt{checkWorkspace} (the type checker, documented in
  \texttt{typing-and-evaluation-rules.tex}), then refreshes \emph{both}
  surfaces from the same diagnostic list, so they cannot disagree. Called
  from the workspace's debounced type-check timer
  (\texttt{app.ts}) and from \texttt{showProblems()}
  (\texttt{app.ts:480--483}). \\
\texttt{\seqsplit{applyBlockWarnings(workspace, diags)}} &
  \texttt{typeDiagnostics.ts:23--34} --- surface \#1: a warning icon
  \emph{on the block itself}, via \texttt{block.setWarningText(...,
  WARNING\_OWNER\_ID)}. Tracks \texttt{markedBlockIds} (l.11) across calls
  so a block whose diagnostic disappeared gets its warning explicitly
  cleared (l.25--29) rather than accumulating stale icons. \\
\texttt{\seqsplit{renderProblemsPanel(workspace, diags)}} &
  \texttt{typeDiagnostics.ts:76--96} --- surface \#2: the Problems list
  itself (\texttt{\#bottom-problems-list}) plus two badge counters
  (\texttt{\#bottom-problems-count} on the tab, \texttt{\#status-problems-count}
  in the status bar (see \texttt{ui-design-and-layout.tex}), both toggling
  a \texttt{has-errors} class when at least one
  diagnostic is an error rather than a warning. \\
\texttt{\seqsplit{renderProblemList(list, workspace, diags)}} &
  \texttt{typeDiagnostics.ts:46--74} --- builds one
  \texttt{button.problem-item} per diagnostic, sorted errors-first
  (l.56--58); each button's click handler is exactly
  \texttt{locateBlock(workspace, diag.blockId)}. \\
\texttt{\seqsplit{locateBlock(workspace, blockId)}} &
  \texttt{typeDiagnostics.ts:37--44} --- \textbf{the canonical
  click-to-locate implementation}, reused verbatim by the Typing tab
  (\texttt{typingPanel.ts:26--30}, per \texttt{typing-code-outline-tabs.tex}):
  \texttt{workspace.getBlockById(blockId)}, then
  \texttt{workspace.centerOnBlock(blockId)}, then
  \texttt{Blockly.common.setSelected(block)} (which, unlike
  \texttt{BlockSvg.select()}, first clears whatever was previously
  selected), and finally dispatches a \texttt{'bmj:problem-located'}
  \texttt{CustomEvent} on \texttt{window} so other listeners (if any) can
  react to a problem having been jumped to. \\
\bottomrule
\end{longtable}
\end{center}

\subsection{Live highlight: CESK and A vs B}

The two machine steppers extend the same idea from "click once to
locate" to "continuously highlight the machine's current focus, and make
\emph{every} rendered value/frame/kontinuation entry individually
clickable." Both are documented in full implementation depth in
\texttt{cesk-machine-and-stepper.tex}; this section gives only the sync
mechanism, for completeness of the bottom-panel inventory.

\begin{center}
\begin{longtable}{@{}p{3.4cm}p{4.4cm}p{6.8cm}@{}}
\toprule
\textbf{Tab} & \textbf{Function} & \textbf{What it does} \\
\midrule
\endhead
CESK &
  \texttt{\seqsplit{locateProvenance(blockId)}} ---
  \texttt{\seqsplit{stepperPanel.ts:90--100}} &
  Click-to-locate on any \texttt{.stepper-provenance} element (control
  text, a frame's method name, a heap box's title, a kontinuation entry) ---
  identical three-call idiom to \texttt{locateBlock} above. \\
CESK &
  \texttt{\seqsplit{setHighlight(blockId)}} ---
  \texttt{\seqsplit{stepperPanel.ts:119--133}} &
  \emph{Continuous} highlight: called from \texttt{renderAll()}
  (l.488) after every step with \texttt{current.focusBlockId}, so the
  program surface always shows exactly where the machine currently is,
  independent of any click. Un-highlights the previous focus block first
  (l.122--124). \\
A vs B &
  \texttt{\seqsplit{setHighlights(blockIds)}} ---
  \texttt{\seqsplit{comparePanel.ts:71--83}} &
  The two-machine analogue: both models' current focus blocks are
  highlighted \emph{simultaneously} on the one shared main workspace
  (\texttt{getWorkspace()} at l.72 is the same closure passed to
  \texttt{stepperPanel}), so a user watching the lockstep view sees both
  machines' positions in the same program surface at once. \\
\bottomrule
\end{longtable}
\end{center}

Both stepper panels take a \texttt{GetWorkspace} closure returning the
\emph{live} main \texttt{Blockly.WorkspaceSvg} at \texttt{init} time
(\texttt{stepperPanel.ts:857}, \texttt{comparePanel.ts}'s equivalent) ---
this is what makes "live highlight" possible at all: the machine
(\texttt{injectMachine(workspace, model)}) is built by reading the real
workspace's blocks, so the \texttt{Blockly.Block} references a
\texttt{MachineState}'s control/frames/heap carry (as \texttt{blockId}
strings, resolved back via \texttt{getBlockById}) always resolve in that
same live workspace.

\subsection{Independent copy: Call-by-Structure and Call-by-Value}
\label{sec:cbs-cbv}

These two tabs are invoked differently from everything else in the dock:
not from a tab click, but from a \textbf{right-click context menu} on a
specific method-call block in the main workspace.

\begin{center}
\begin{longtable}{@{}p{4.0cm}p{10.2cm}@{}}
\toprule
\textbf{Symbol} & \textbf{Role} \\
\midrule
\endhead
\texttt{\seqsplit{registerMiniJavaContextMenus()}} &
  \texttt{contextMenus.ts:17--62} --- registers two
  \texttt{Blockly.ContextMenuRegistry} items, \texttt{miniJavaVizStructure}
  and \texttt{miniJavaVizValue}, scoped to \texttt{ScopeType.BLOCK} and
  shown only when the right-clicked block is \texttt{mj\_expr\_method\_call}
  (\texttt{contextMenus.ts:28,38}); each callback calls
  \texttt{openVisualization(kind, scope.block)}. \\
\texttt{\seqsplit{openVisualization(kind, block)}} &
  \texttt{visualizationPanel.ts:192--197} --- stores the seed
  \texttt{block} in \texttt{views[kind].block}, activates that tab, opens
  the dock, and calls \texttt{renderView(kind)}. \\
\texttt{renderView(kind)} &
  \texttt{visualizationPanel.ts:175--190} --- clears (or lazily creates via
  \texttt{ensureWorkspace}, l.95--98) a \emph{brand-new}
  \texttt{Blockly.WorkspaceSvg} injected into that tab's own
  \texttt{.viz-host} (\texttt{injectVisualizationWorkspace}, l.78--93,
  with its own zoom/grid/theme settings), then calls
  \texttt{renderMiniJavaReduction} or \texttt{renderCopiedReductionRoot}
  (\texttt{core/semantics/minijavaReduction.ts}) to populate it with a
  \emph{copy} of the reduction tree rooted at the seed block. \\
\bottomrule
\end{longtable}
\end{center}

\textbf{Why no click-to-locate back to the main workspace:} the blocks
inside the Structure/Value workspace are independent copies (fresh block
ids, via \texttt{Blockly.serialization.blocks}-style duplication inside
\texttt{minijavaReduction.ts}) --- this is a direct visual consequence of
the substitution semantics these tabs illustrate (\texttt{typing-and-
evaluation-rules.tex}, \S4.3): each reduction step is rendered as its own
\emph{self-contained} block tree, and two occurrences of a substituted
value are, by the structure-preservation invariant, never the same node.
There is accordingly no single "the" main-workspace block a click inside
this copy could unambiguously jump back to. \texttt{disposeVizWorkspaces()}
(\texttt{visualizationPanel.ts:199--214}) tears down and, if that tab was
open, re-renders these private workspaces on a theme switch, confirming
they are wholly owned by the dock and carry no reference back into the
main workspace's block set.

\subsection{Independent copy: Rewrite}

The Rewrite (substitution stepper) tab, \texttt{src/core/ui/substPanel.ts}
(295 lines), follows the same "own private workspace" pattern as
\S\ref{sec:cbs-cbv}, but is entered from the tab strip directly (like CESK)
rather than from a context menu, and highlights within that private
workspace continuously as it steps.

\begin{center}
\begin{longtable}{@{}p{4.0cm}p{10.2cm}@{}}
\toprule
\textbf{Symbol} & \textbf{Role} \\
\midrule
\endhead
\texttt{getWorkspace} (main) &
  \texttt{substPanel.ts:32} --- used \emph{only} to seed the initial tree
  (from the \texttt{println} argument in \texttt{main}, per
  \texttt{typing-and-evaluation-rules.tex} \S4.3.1) and to detect
  staleness on a live edit (\texttt{attachWorkspaceListener},
  \texttt{substPanel.ts:97--118} --- same \texttt{BLOCK\_CREATE/DELETE/
  CHANGE/MOVE} listener shape as every other stepper). It is
  \textbf{never} the target of a click-to-locate. \\
own highlight (redex) &
  \texttt{substPanel.ts:176} --- \texttt{\seqsplit{host.getBlockById(current.redexId).setHighlighted(true)}},
  where \texttt{host} is the tab's own private workspace (populated by
  re-rendering \texttt{current.tree} --- a \texttt{TreeNode}, not a
  \texttt{Blockly.Block} --- into fresh blocks each step). The highlighted
  block is always inside \emph{that} tree, never the main program's. \\
\bottomrule
\end{longtable}
\end{center}

This is the clearest illustration of the regime split in this note: CESK
and A~vs~B evaluate the \emph{actual program blocks by reference}
(\texttt{Blockly.Block} objects looked up live by id), so highlighting them
\emph{is} highlighting the main workspace; Structure/Value/Rewrite evaluate
a \emph{copied, self-contained syntax tree} (the substitution semantics'
own state representation, per the companion evaluation-rules note), so
their highlight necessarily targets their own private copy.

\subsection{Text only: Output}

\texttt{src/core/ui/programConsole.ts} (50 lines) is the simplest tool in
the dock and syncs with nothing block-related at all --- its "source of
truth" is a plain string log, not a workspace.

\begin{center}
\begin{longtable}{@{}p{4.0cm}p{10.2cm}@{}}
\toprule
\textbf{Function} & \textbf{Role} \\
\midrule
\endhead
\texttt{\seqsplit{mirrorProgramOutput(source, lines, note)}} &
  \texttt{programConsole.ts:24--40} --- overwrites
  \texttt{\#bottom-program-output} with \texttt{[source]} plus the joined
  \texttt{println} lines, flashing \texttt{is-console-changed} when the
  body actually changed (l.31--34). \textbf{Whichever tool produced output
  most recently owns the console}: \textbf{Run} (\texttt{app.ts:340--351}),
  the CESK stepper (\texttt{stepperPanel.ts:477}), and A~vs~B all call
  this same function, so the Output tab is a shared monitor, not one
  tool's private log. \\
\texttt{\seqsplit{appendConsoleLog(message)}} &
  \texttt{programConsole.ts:43--50} --- appends one line (used for
  editor-action notices like export/clipboard) without replacing the
  mirrored run output. \\
\bottomrule
\end{longtable}
\end{center}

% =============================================================================
\section{Full tab inventory}

\begin{center}
\begin{longtable}{@{}p{2.4cm}p{2.8cm}p{3.2cm}p{6.0cm}@{}}
\toprule
\textbf{Tab} & \textbf{\texttt{data-kind}} & \textbf{Sync regime} & \textbf{Owning module} \\
\midrule
\endhead
Problems & \texttt{problems} & Live highlight & \texttt{typeDiagnostics.ts} \\
Output & \texttt{output} & Text only & \texttt{programConsole.ts} \\
Call-by-Structure & \texttt{structure} & Independent copy & \texttt{visualizationPanel.ts} + \texttt{minijavaReduction.ts} \\
Call-by-Value & \texttt{value} & Independent copy & \texttt{visualizationPanel.ts} + \texttt{minijavaReduction.ts} \\
CESK & \texttt{machine} & Live highlight & \texttt{stepperPanel.ts} \\
A vs B & \texttt{compare} & Live highlight & \texttt{comparePanel.ts} \\
Rewrite & \texttt{subst} & Independent copy & \texttt{substPanel.ts} \\
\bottomrule
\end{longtable}
\end{center}

% =============================================================================
\section{HTML wiring}

\begin{center}
\begin{longtable}{@{}p{3.6cm}p{10.6cm}@{}}
\toprule
\textbf{Element} & \textbf{Location} \\
\midrule
\endhead
\texttt{\#viz-dock} &
  \texttt{index.html:243} --- the dock root; \texttt{data-open} toggled by
  \texttt{setVizOpen}. \\
\texttt{.viz-dock-tabs} &
  \texttt{index.html:247--268} --- the seven \texttt{button.viz-tab[data-kind]}
  elements, in \texttt{ALL\_KINDS} order. \\
\texttt{\#viz-dock-info} &
  \texttt{index.html:270} --- the per-tab description line
  (\S\ref{sec:info}). \\
\texttt{.viz-dock-tools} &
  \texttt{index.html:271--276} --- \textbf{Re-run}
  (\texttt{\#viz-rerun}), \textbf{Arrange} (\texttt{\#viz-arrange}),
  \textbf{Maximize} (\texttt{\#viz-maximize}), \textbf{Hide}
  (\texttt{\#viz-collapse}). \\
\texttt{\seqsplit{.viz-host[data-kind]}} &
  \texttt{index.html:279--328} --- seven host
  containers, one per tab; \texttt{problems}/\texttt{output} are
  \texttt{.bottom-passive-host} (plain content, no injected Blockly
  workspace); \texttt{structure}/\texttt{value}/\texttt{machine}/\texttt{compare}/\texttt{subst}
  each host either an injected reduction workspace or the CESK/Compare/Rewrite
  stepper layout (documented in \texttt{cesk-machine-and-stepper.tex} and
  \texttt{typing-and-evaluation-rules.tex} respectively). \\
\texttt{\#viz-resizer} &
  \texttt{index.html:244} --- the drag handle for dock height
  (\S1, \texttt{initResizer}). \\
\bottomrule
\end{longtable}
\end{center}

% =============================================================================
\section*{Note on scope}
\addcontentsline{toc}{section}{Note on scope}
This note covers the bottom panel's \emph{shell} (tab switching, open/close,
resize) and its \emph{synchronization} with the workspace. The CESK and A~vs~B
tabs' machine semantics and internal rendering (Control/Environment/Store/Kontinuation
columns, SVG heap arrows, garbage collection) are covered in
\texttt{cesk-machine-and-stepper.tex} and \texttt{garbage-collection.tex}. The
Rewrite tab's substitution semantics are covered in
\texttt{typing-and-evaluation-rules.tex}. The three inspector-column tabs
(Editable Code, Typing, Outline) are a separate panel, covered in
\texttt{typing-code-outline-tabs.tex}. The workbench shell the dock sits
inside (activity bar, perspectives, responsive layout) is covered in
\texttt{ui-design-and-layout.tex}.

\end{document}
